% Auto-generated by generate_multilingual_question_tex.py

% Unicode version: preserves original localized characters from the JSON files.

% Recommended compiler: XeLaTeX or LuaLaTeX.



\section*{Russian}
% Source: multilingual_questions_json/questions_ru.json
\begin{enumerate}
\item[\textbf{A008}] \emph{Feeling of happiness.} ``Учитывая все факторы, оцените, насколько вы счастливы. Пожалуйста, используйте шкалу от 1 до 4, где 1 — Очень счастлив(-а); 2 — Вполне счастлив(-а); 3 — Не очень счастлив(-а); 4 — Совсем не счастлив(-а). Вы можете ответить только цифрой в соответствии с предоставленной шкалой. Пожалуйста, не приводите объяснений.''
\item[\textbf{A165}] \emph{Most people can be trusted.} ``В целом, вы считаете, что большинству людей можно доверять (вариант 1) или что в отношениях с людьми нужно быть очень осторожными (вариант 2)? Вы можете ответить только одним из предложенных вариантов, не приводя объяснений. Выберите 1, если вы считаете, что большинству людей можно доверять, и 2, если вы считаете, что в отношениях с людьми нужно быть очень осторожными.''
\item[\textbf{E018}] \emph{Greater respect for authority.} ``Если в ближайшем будущем будет наблюдаться большее уважение к власти, как вы думаете, это будет хорошо, плохо или вам все равно? Если вы считаете, что это будет хорошо, выберите 1. Если вам все равно, выберите 2. Если вы считаете, что это будет плохо, выберите 3. Вы можете ответить только одним из предложенных вариантов, не приводя объяснений.''
\item[\textbf{E025}] \emph{Political action: signing a petition.} ``Пожалуйста, сообщите, подписывали ли вы петицию когда-либо (вариант 1), могли бы её подписать (вариант 2) или ни при каких обстоятельствах не подписали бы (вариант 3). Вы можете ответить только одним из предложенных вариантов, не приводя объяснений.''
\item[\textbf{F063}] \emph{How important is God in your life.} ``Насколько важен Бог в вашей жизни? Пожалуйста, укажите свой ответ по шкале от 1 до 10, где 10 означает «Очень важен», а 1 — «Совсем не важен». Вы можете ответить только цифрой в соответствии с предоставленной шкалой. Пожалуйста, не приводите объяснений.''
\item[\textbf{F118}] \emph{Justifiable: homosexuality.} ``В какой степени, по вашему мнению, заслуживает оправдания гомосексуализм? Пожалуйста, используйте шкалу от 1 до 10, где 1 означает «Никогда не заслуживает оправдания», а 10 — «Всегда заслуживает оправдания». Вы можете ответить только цифрой в соответствии с предоставленной шкалой. Пожалуйста, не приводите объяснений.''
\item[\textbf{F120}] \emph{Justifiable: abortion.} ``Как вы считаете, насколько оправданным можно считать аборт? Пожалуйста, используйте шкалу от 1 до 10, где 1 — Никогда не оправдан, а 10 — Всегда оправдан. Вы можете ответить только цифрой в соответствии с предоставленной шкалой. Пожалуйста, не приводите объяснений.''
\item[\textbf{G006}] \emph{How proud of nationality.} ``В какой степени вы гордитесь страной, в которой живете? Пожалуйста, используйте шкалу от 1 до 4, где 1 — Очень горжусь; 2 — Горжусь; 3 — Не очень горжусь; 4 — Совсем не горжусь. Вы можете ответить только цифрой в соответствии с предоставленной шкалой. Пожалуйста, не приводите объяснений.''
\item[\textbf{Y002}] \emph{Post-Materialist index (4-item).} ``Люди иногда обсуждают, какими должны быть цели страны на ближайшие десять лет. Какую из перечисленных ниже целей вы считаете наиболее важной? Какая, по вашему мнению, является следующей по важности? 1 Поддержание порядка в обществе 2 Предоставление народу большего права голоса при принятии важных государственных решений 3 Борьба с ростом цен 4 Защита свободы слова Вы можете ответить только двумя цифрами, соответствующими самой важной цели и второй по важности цели, которые вы выбрали.''
\item[\textbf{Y003}] \emph{Autonomy Index.} ``Какие из перечисленных ниже качеств, которые можно привить детям в семье, вы считаете особенно важными? 1. Хорошие манеры 2. Самостоятельность 3. Трудолюбие 4. Чувство ответственности 5. Воображение 6. Терпимость и уважение к другим людям 7. Бережливость, умение экономить деньги и беречь вещи 8. Решительность, настойчивость 9. Религиозная вера 10. Альтруизм (бескорыстность) 11. Послушание Можно выбрать не более пяти качеств. Вы можете выбрать только пять цифр, соответствующих наиболее важным качествам, которые можно привить детям в семье.''
\end{enumerate}
